Lu matches upper hits on the descent and obtains the first-order
correction corresponding to $U(a)-f(a)\Tr(t)/f(t)$
\cite[supplement, p.~16, (S5.12)--(S5.13)]{LuDegree}.
Two accuracies are needed here. Before the reference hits agree, the
exponential rate locates each upper deficit. After matching, the two
terms in the derivative of the lower-intercept difference have a
common exponential rate through the crest. Their first-order log
ratio determines the sign away from the crest.

\begin{lemma}[Single-half deficits and a common exponent across a crest]\label{lem:common-crest}
Use the scaled legs of Lemma~\ref{lem:variational}, with potential
coordinate $u=r^2/2$ and half-profiles $p_\sigma(r)=\phi_\sigma(u)$,
$\sigma=\pm1$. Let $t_{b,\sigma}\to t_b$ be later reference radii and
fix a finite $T_{\max}>t_b^2/2$. On $\mathcal I=[0,T_{\max}]$ set
$\widetilde E(u)=E(\sqrt{2u})$ and $\widetilde Q(u)=Q(\sqrt{2u})$.
Assume the profiles vanish at the axis, $E(0)=Q(0)=0$, and, as $\delta\downarrow0$,
\[
 \|\phi_\sigma-\widetilde E-\sigma\delta\widetilde Q\|_{C^1(\mathcal I)}
     =o(\delta),\qquad
 \sup_{\delta,\sigma}\operatorname{Var}_{\mathcal I}(\phi_\sigma')<\infty,
 \qquad \sqrt\eps+\eps|\log\eps|=o(\delta).
\]
On each compact subinterval of $(0,T_{\max}]$ assume
$\phi_\sigma\to\widetilde E$ in $C^2$.
For $0<r\le\sqrt{2T_{\max}}$, the even $E$ rises to a simple maximum
at $1$ and then decreases through zero at $1<g<t_b$, with nonzero
slope off the crest. The perturbed profiles likewise have a regular
ascent and descent outside an initial interval $0\le u\le\eta_\delta$,
where $\eta_\delta\to0$ and
$\int_0^{\eta_\delta}|\phi_\sigma'|^2\dd u=o(\delta)$.
Assume transverse axis hits for the turns used below, including the turns
up to each reference endpoint $T_{b,\sigma}=t_{b,\sigma}^2/2$.
Put $f=E'/r$, $U(t)=\int_0^t(fQ'-f'Q)\dd r$, and $\Tr(t)=U(t)-U(1)$.
\begin{enumerate}
  \item[\textup{(i)}] \emph{Single-half deficit rate.} Without assuming
  equality of the reference upper hits, write
  $\mathcal A_{\mathrm{out},\sigma}(S)=\int_0^S \phi_\sigma'\dd m_{\sigma,S}$.
  Uniformly on compact subsets of $T\in[0,t_b^2/2)$ bounded away from $T_{b,\sigma}$,
  \[
    -\eps\log\bigl(X_\sigma(T_{b,\sigma})-X_\sigma(T)\bigr)
    = \min_{S\in[T,T_{b,\sigma}]} \mathcal A_{\mathrm{out},\sigma}(S)
      + O\bigl(\sqrt\eps+\eps|\log\eps|\bigr).
  \]
  In particular, for a descending turn $T$ bounded away from the perturbed
  crest $r_{\mathrm{crest},\sigma}$ and from $T_{b,\sigma}$, the minimum is
  attained at $S=T$ and equals
  $\int_{r_{\mathrm{crest},\sigma}}^{\sqrt{2T}}(\partial_r p_\sigma)^2\dd r/r$.
  \item[\textup{(ii)}] \emph{Matched reference upper hits.} Assume in addition
  that $X_+(T_{b,+})=X_-(T_{b,-})$. Let $t_-(t)$ match the upper hit of the
  plus turn $t$, write $\Delta_Y(t)=Y_+(t^2/2)-Y_-(t_-^2/2)$ and
  $J_\pm=J_{i,\pm}+J_{o,\pm}$ at these actual turns, and for $0<t<g$ let $a=t$
  on the ascent and let $a\le 1$ satisfy $E(a)=E(t)$ on the descent. On compact
  subintervals of $(0,g)$, including those containing $1$, both positive terms
  \[
   \mathcal L=e^{-J_{i,+}},\qquad \mathcal N=e^{J_{o,+}-J_{o,-}-J_{i,-}},\qquad
   \Delta_Y'=\eps t(\mathcal L-\mathcal N)
  \]
  satisfy $-\eps\log \mathcal L,\,-\eps\log \mathcal N\to W(t)=\int_0^a f(r)^2r\dd r$
  uniformly. On compact sets of regular turns the finer limit is
  \begin{equation}\label{eq:crest-splitting}
   \frac{\eps}{2\delta}\log(\mathcal N/\mathcal L)\longrightarrow
   \begin{cases}U(t),&0<t<1,\\
   U(a)-f(a)\Tr(t)/f(t),&1<t<g .
   \end{cases}
  \end{equation}
  If in addition $\Delta_Y(a_*)=0$ for $0<a_*<1$, the signs of the ordinary
  displacement at $a_i<a_*$ and at its descending partner
  $z_i>1$ with $E(z_i)=E(a_i)$ are the negative and the positive,
  respectively, of the corresponding nonzero expressions in
  \eqref{eq:crest-splitting}. A nonzero $U(a_*)$ makes the marked orbit simple.
\end{enumerate}
\end{lemma}

\begin{proof}
For part~\textup{(i)}, suppress $\sigma$ and write $T_b=T_{b,\sigma}$.
For each $0<S\le T_b$ the two envelope actions are
\[
 \mathcal A_{\rm out}(S)=\int_0^S\phi'\dd m_S,\qquad
 I(S)=\int_0^S\phi'\dd l_S,
\]
where $m_S,l_S$ are the suffix extrema in \eqref{eq:running-extrema}.
Integration by parts on $[0,S]$, using \eqref{eq:turn-derivs} and
\eqref{eq:envelope-bounds}, bounds the errors in $-\eps J_o(S)$ and
$\eps J_i(S)$ by $\sqrt{2\eps S}\,(2\|\phi'\|_{\infty,\mathcal I}+
\operatorname{Var}_{\mathcal I}\phi')$. Thus one bound applies from the axis to
either reference turn for all sufficiently small $\delta$.
The envelopes change by at most $\|\phi'\|_\infty |T-T'|$ on their
common interval. The actions are uniformly Lipschitz in the turn by
the same integration by parts. They converge uniformly to those of $E$:
the envelopes converge in sup norm, their variations are bounded, and
the limiting slope is of bounded variation. The limiting upper action
is zero up to the crest and then equals $\int_1^t f^2 r\dd r$, strictly
increasing. The lower action for $0<t<g$ is $W(t)$.

Put $X_*=X_\sigma(T_b)$.
The exact deficit $X_*-X_\sigma(T)=\int_T^{T_b}\eps e^{J_{o,\sigma}(S)}\dd S$
has negative scaled logarithm within
$O(\sqrt\eps+\eps|\log\eps|)$ of $\min_{[T,T_b]}\mathcal A_{\rm out}(S)$,
uniformly away from $T_b$. For the upper bound on the integral use the minimum cost on the
whole interval. For the lower bound integrate on a segment of length
$\eps$ adjacent to a minimizer, where Lipschitz continuity bounds the
increase in action by $O(\eps)$. At a regular descending turn the
suffix maximum follows the descent from the perturbed crest, so the
minimum outgoing cost occurs at the first endpoint and equals
$\int_{r_{\mathrm{crest},\sigma}}^{\sqrt{2T}}(\partial_r p_\sigma)^2\dd r/r$.
This proves part~\textup{(i)}.

For part~\textup{(ii)}, differentiate the equality of the upper hits
$X_+(t^2/2)=X_-(t_-(t)^2/2)$ using \eqref{eq:turn-derivs} to obtain the
formula for $\Delta_Y'$. Part~\textup{(i)} distinguishes the branches for
paired upper hits. On the ascent the limiting upper
height increases to $E(1)$, which rules out a partner beyond the crest
for a regular rising turn. On the descent the limiting deficit cost is
strictly increasing and positive, whereas its value is zero on the ascent.
For a compact test range in $(0,g)$ choose a fixed descending radius
between $g$ and $t_b$. Its limiting deficit cost exceeds all the test
costs, so monotonicity of $X_-$ keeps the partner below this radius.
As $t\to1$, a limiting partner must attain both upper height $E(1)$
and zero deficit cost. A radius below $1$ misses the former and a
radius above $1$ misses the latter. Hence $t_-(t)-t\to0$ uniformly
through the crest as well. The costs of $\mathcal L,\mathcal N$ are therefore
$I_+$ and $I_-+\mathcal A_{{\rm out},+}-\mathcal A_{{\rm out},-}$ up to $O(\sqrt\eps)$, both with limit $W$.

Write $\kappa_+=p_+(t)$ and $\kappa_-=p_-(t_-(t))$ for the
actual scaled turning ordinates of the two arcs with common upper hit.
For a fixed level $\kappa$ on the rising branch, the endpoint defined
by $E(a)+\delta Q(a)=\kappa$ has derivative $-Q(a)/E'(a)$. In the variation of
$\int (E'+\delta Q')^2\dd r/r$ its boundary term cancels $f(a)Q(a)$
in $2\int fQ'$. This gives $U(a)$ up to the initial endpoint term,
which vanishes here ($Q(0)=0$ for the half-profiles).
The analogous descending integral from the crest gives $\Tr(t)$, since
$f(1)=0$. On the initial interval a suffix envelope follows the profile or is
constant, so its action there is bounded by
$\int_0^{\eta_\delta}|\phi_\sigma'|^2\dd u=o(\delta)$.
The $C^1(\mathcal I)$ expansion gives the same bound for the contribution of
$(\widetilde E')^2$ there. Local $C^2$ convergence at the simple crest
controls the perturbed maximum, while the first variations use the
$C^1(\mathcal I)$ remainder.
On the ascent the matched terminal levels $\kappa_\pm$ agree to
$O(\sqrt\eps)=o(\delta)$, so the inward variation gives the first
line of \eqref{eq:crest-splitting}.

On a compact interval of regular descending levels write
$O_\sigma(k)$ for the outgoing envelope action at the turn of level $k$.
Write $z(k)>1$ for the descending inverse of $E$. At zero perturbation
$O_0(k)=\int_1^{z(k)} f(r)^2r\dd r$, hence $O_0'(k)=f(z(k))<0$.
The fixed-level variation just computed gives
$O_\sigma(k)=O_0(k)+\sigma\delta\Tr(z(k))+o(\delta)$ uniformly there.
At the paired turns $T_+=t^2/2$, $T_-=t_-(t)^2/2$, both the upper
hit and the later reference hit agree. Hence the two actual deficits
are the same positive number
\[
 \mathfrak d_*=X_+(T_{b,+})-X_+(T_+)=X_-(T_{b,-})-X_-(T_-).
\]
On a compact set of regular descending turns the pairs stay a fixed
distance from their respective reference endpoints. The outgoing
envelope action on $[T_\sigma,T_{b,\sigma}]$ is minimized at
$T_\sigma$. With $e_\eps=\sqrt\eps+\eps|\log\eps|$, the single-half
deficit estimate above thus gives
$O_\sigma(\kappa_\sigma)=-\eps\log\mathfrak d_*+O(e_\eps)$ for each sign,
uniformly on this set. The bound for the actual outgoing leg now gives
\[
 -\eps\bigl(J_{o,+}(T_+)-J_{o,-}(T_-)\bigr)
 = O_+(\kappa_+)-O_-(\kappa_-)+O(\sqrt\eps)
 = O(e_\eps)=o(\delta).
\]
Since $|O_0'|$ is bounded away from zero, the mean value theorem first
gives $\kappa_+-\kappa_-=O(\delta)$. Both levels tend to $E(t)$, so
applying the same expansion at the two moving levels gives
\[
 o(\delta)=O_+(\kappa_+)-O_-(\kappa_-)
 =f(t)(\kappa_+-\kappa_-)+2\delta\Tr(t)+o(\delta).
\]
Thus $\eps\log(\mathcal N/\mathcal L)=\eps(J_{i,+}-J_{i,-})+o(\delta)$.
The inward action retains the change of level through its derivative
$f(a)$, in addition to its fixed-level variation $U(a)$. Consequently
\[
 \frac{\eps}{2\delta}\log\frac{\mathcal N}{\mathcal L}
 =U(a)+f(a)\frac{\kappa_+-\kappa_-}{2\delta}+o(1)
 =U(a)-f(a)\frac{\Tr(t)}{f(t)}+o(1).
\]
All remainders in this first-order calculation are uniform on compact
sets of regular turns.

On $[a_i,a_*]$ the cost $W$ increases strictly. On $[a_*,z_i]$
it first increases to $W(1)$ and then decreases strictly to
$W(z_i)=W(a_i)<W(a_*)$, since $E(z_i)=E(a_i)<E(a_*)$.
Each test is therefore the unique minimum on its interval to the mark. In a fixed neighborhood of that endpoint the larger of
$\mathcal L,\mathcal N$ exceeds the smaller by $\exp(c\delta/\eps)$ for some $c>0$.
Outside that neighborhood the limiting cost is at least
$W(\text{test})+3\gamma$ for some $\gamma>0$. On a shorter interval
at the test it is below $W(\text{test})+\gamma$.
Uniform logarithmic convergence makes the latter signed mass dominate.
Integrating in the appropriate direction and using
$\sgn\oscD=-\sgn \Delta_Y$ proves the assertions. At the mark the multiplier
is $e^{J_--J_+}$.
\end{proof}
